\chapter{Appendix}
\section{Continued Fractions over Finite Fields}\label{sec:cont-frac-fq}

Let $r_0,r_1\in\bbF_q[X]$ with $r_1\neq 0$. Applying the Euclidean algorithm yields a sequence
$$r_i=A_ir_{i+1}+r_{i+2},\quad i=0,1,\dots,s,$$
where $0\leq\deg(r_{i+1})<\deg(r_i)$ for $i=1,\dots,s$ and $r_{s+2}=0$. Here $A_0,A_1,\dots,A_s\in\bbF_q[X]$, with $A_1,\dots,A_s$ of positive degree. This gives the \emph{continued fraction expansion}
$$\frac{r_0}{r_1}=A_0+\dfrac{1}{A_1+\dfrac{1}{A_2+\cdots}}=:[A_0,A_1,A_2,\dots,A_s].$$
Define polynomials $P_i,Q_i\in\bbF_q[X]$ for $i=-1,0,1,\dots,s$ by the recurrences
\begin{align*}
  P_{-1}=1,\quad P_0=A_0,\quad & P_i=A_iP_{i-1}+P_{i-2}\quad (i=1,\dots,s), \\
  Q_{-1}=0,\quad Q_0=1,\quad   & Q_i=A_iQ_{i-1}+Q_{i-2}\quad (i=1,\dots,s).
\end{align*}
It is clear that $\deg(P_{i-1})<\deg(P_i)$ and $\deg(Q_{i-1})<\deg(Q_i)$ for $i=1,\dots,s$. We extend the definition of degree by $\deg(f/g)=\deg(f)-\deg(g)$ for a rational function $f/g$.

\begin{lemma}{Rational Approximants via Convergents}{lem:cf-rational-approx}
  For any rational function $\rho$ of nonnegative degree,
  $$[A_0,A_1,\dots,A_{i-1},\rho]=\frac{\rho P_{i-1}+P_{i-2}}{\rho Q_{i-1}+Q_{i-2}}\quad\text{for }i=1,\dots,s+1.$$
\end{lemma}
\begin{proof}
  We proceed by induction on $i$. For $i=1$ both sides equal $A_0+\rho^{-1}=\rho^{-1}(A_0\rho+1)=(\rho P_0+P_{-1})/(\rho Q_0+Q_{-1})$. Suppose the identity holds for some $i$, $1\leq i\leq s$. Since $A_i+\rho^{-1}$ has positive degree, applying the inductive hypothesis gives
  $$[A_0,\dots,A_i,\rho]=\frac{(A_i+\rho^{-1})P_{i-1}+P_{i-2}}{(A_i+\rho^{-1})Q_{i-1}+Q_{i-2}}=\frac{\rho P_i+P_{i-1}}{\rho Q_i+Q_{i-1}},$$
  using the recurrence relations. \qed
\end{proof}

\begin{corolary}{Convergents Formula}{cor:piqi-is-expression}
  For $i=0,1,\dots,s$ we have $[A_0,A_1,\dots,A_i]=P_i/Q_i$.
\end{corolary}

\begin{lemma}{Representation of $r_0/r_1$ via Convergents}{lem:cf-r0r1-rep}
  For $i=0,1,\dots,s$ we have
  $$\frac{r_0}{r_1}=\frac{P_i+\beta_i P_{i-1}}{Q_i+\beta_i Q_{i-1}},$$
  where $\beta_i=r_{i+2}/r_{i+1}$ is a rational function of negative degree.
\end{lemma}
\begin{proof}
  We proceed by induction. For $i=0$, $(P_0+\beta_0P_{-1})/(Q_0+\beta_0Q_{-1})=(A_0+r_2/r_1)/(1)=(A_0r_1+r_2)/r_1=r_0/r_1$. The inductive step follows from the recurrences and $\beta_i^{-1}=A_{i+1}+\beta_{i+1}$. \qed
\end{proof}

\begin{lemma}{Determinant Identity for Convergents}{lem:cf-determinant}
  For $i=0,1,\dots,s$ we have $P_iQ_{i-1}-P_{i-1}Q_i=(-1)^{i-1}$.
\end{lemma}
\begin{proof}
  Induction. For $i=0$: $P_0Q_{-1}-P_{-1}Q_0=0\cdot 0-1\cdot 1=-1=(-1)^{-1}$. Assuming the identity at $i$,
  $$P_{i+1}Q_i-P_iQ_{i+1}=(A_{i+1}P_i+P_{i-1})Q_i-P_i(A_{i+1}Q_i+Q_{i-1})=-(P_iQ_{i-1}-P_{i-1}Q_i)=(-1)^i. \qed$$
\end{proof}

\begin{corolary}{Coprimality of Convergents}{cor:piqi-coprime}
  For $i=0,1,\dots,s$ we have $\gcd(P_i,Q_i)=1$.
\end{corolary}
\begin{proof}
  If $d_i=\gcd(P_i,Q_i)$ then \lemref{lem:cf-determinant} shows $d_i$ divides $(-1)^{i-1}$, so $d_i=1$. \qed
\end{proof}

\section{A Product Identity for Roots of Unity}

The following identity, used in the factorisation of lifted $L$-functions, expresses a product of linear factors twisted by roots of unity as a perfect power.

\begin{lemma}{Roots-of-Unity Product Identity}{lem:roots-unity-product}
  Let $n,m\in\bbN$ and set $d=\gcd(m,n)$. Then
  $$\prod_{t=1}^{n}\lt(1-e^{2\pi i\nicefrac{mt}{n}}\,X\rt)=\lt(1-X^{\nicefrac{n}{d}}\rt)^{d}.$$
\end{lemma}
\begin{proof}
  Let $\omega=e^{2\pi i/n}$ be a primitive $n$-th root of unity, so the factors are $1-\omega^{mt}X$ for $t=1,\dots,n$. Since $t=n$ gives $\omega^{mn}=1=\omega^0$, the product is the same as $\prod_{t=0}^{n-1}(1-\omega^{mt}X)$.

  Because $d=\gcd(m,n)$, the element $\omega^m$ is a primitive $(n/d)$-th root of unity. As $t$ ranges over $\{0,1,\dots,n-1\}$, the values $\omega^{mt}$ cycle through the $(n/d)$-th roots of unity exactly $d$ times each. Hence
  $$\prod_{t=0}^{n-1}\lt(1-\omega^{mt}X\rt)=\lt(\prod_{k=0}^{n/d-1}\lt(1-(\omega^m)^k X\rt)\rt)^d.$$
  Since $\omega^m$ is a primitive $(n/d)$-th root of unity, the inner product equals $1-X^{n/d}$ by the standard factorisation of $1-Y^N = \prod_{k=0}^{N-1}(1-\zeta^k Y)$. Therefore the whole product equals $(1-X^{n/d})^d$. \qed
\end{proof}

\section{Elementary Symmetric Polynomials}
Waring's Formula expresses each power-sum symmetric polynomial $s_k=x_1^k+x_2^k+\cdots+x_n^k$ as an explicit polynomial in the elementary symmetric polynomials $\sigma_1,\dots,\sigma_n$. It is the source of the closed-form reduction of lifted Kloosterman sums $K(\chi^{(s)};a,b)$ to the ground-field sum $K(\chi;a,b)$ in \lemref{lem:kloosterman-reduction}: applying Waring's formula to $\omega_1^s+\omega_2^s$ with $\sigma_1=\omega_1+\omega_2=-K(\chi;a,b)$ and $\sigma_2=\omega_1\omega_2=q$ produces the binomial-coefficient expansion. We state the general identity here for reference.
\begin{Theorem}{Waring's Formula}{thm:waring}%
  \index{Waring's formula}
  Let $x_1,\dots,x_n$ be indeterminates, let $$\sigma_j=\sum_{1\leq i_1<i_2<\cdots<i_j\leq n}x_{i_1}x_{i_2}\cdots x_{i_j}\qquad (j=1,\dots,n)$$ be the elementary symmetric polynomials in $x_1,\dots,x_n$, and let $$s_k=x_1^k+x_2^k+\cdots+x_n^k$$ be the $k^{\text{th}}$ power-sum. Then for every integer $k\geq 1$,
  $$s_k=\sum(-1)^{i_2+i_4+i_6+\cdots}\cdot \frac{(i_1+i_2+\cdots+i_n-1)!\cdot k}{i_1!\cdot i_2!\cdots i_n!}\cdot \sigma_1^{i_1}\sigma_2^{i_2}\cdots \sigma_n^{i_n},$$
  where the summation runs over all $n$-tuples $(i_1,\dots,i_n)$ of non-negative integers satisfying $i_1+2i_2+\cdots+ni_n=k$. Each coefficient of $\sigma_1^{i_1}\sigma_2^{i_2}\cdots \sigma_n^{i_n}$ is an integer.
\end{Theorem}

\begin{Theorem}{}{thm:elm-symmetry-spans}
  Suppose $f\in\bbF_q[X_1,\dots,X_n]$ is a symmetric polynomial. Let $\ovX=(X_1,\dots, X_n)$. Then there exists a polynomial $g\in \bbF_q[Y_1,\dots,Y_n]$ such that $$f(X_1,\dots, X_n)=g(\esym_1(\ovX),\dots,\esym_n(\ovX)).$$

  Moreover, \begin{enumerate}[label=(\roman*)]
    \item if $\deg_{X_i}(f)=d$ for all $i\in[n]$ then $\deg G=d$.
    \item if $\deg f=D$ then for every monomial $Y_1^{d_1}\cdots Y_n^{d_n}$  of $g$ with non-zero coefficient then $d_1+2\cdot d_2+\cdots +n\cdot d_n= D$.
  \end{enumerate}
\end{Theorem}

\section{Hasse Derivative}
\begin{Definition}{Hasse Derivative}{}
  Let $\bbF$ be any field. Let $P\in\bbF[X_1,\dots,X_n]$. Then $\ovi\in \bbZ_0^n$ Hasse derivative of $P$ is denoted by $P^{(\ovi)}(X_1,\dots, X_n)$ or $H^{(\ovi)}(P)$ which is the coefficient of $Z^{i_1}\cdots Z_n^{i_n}\eqqcolon \ovZ^{\ovi}$ in the polynomial  $P(X_1+Z_1,\dots, X_n+Z_n)$. Thus $$P(X_1+Z_1,\dots, X_n+Z_n)= \sum_{\ovi} P^{(\ovi)}(X_1,\dots, X_n)\cdot \ovZ^{ \ovi}$$
\end{Definition}

For univariate polynomials we immediately get a closed form formula of $k^{th}$ hasse derivative. Let $f\in\bbF[X]$ where $f(X)=\sum_{i=0}^da_iX^i$. Then $$f^{(k)}(X)=\sum_{i=0}^d \binom{k}{i}\cdot a_iX^{i-k}$$

From the definition we have the following observation
\begin{observation}
  For any $f,f_1,f_2\in\bbF_q[X_1,\dots, X_n]$ and $\alpha\in\bbF$ and for any $\ovk\in\bbZ^n$ we have\begin{enumerate}[label=(\roman*)]
    \item $(f_1+f_2)^{(\ovk)}(X_1,\dots, X_n)=f_1^{(\ovk)}(X_1,\dots, X_n)+f_2^{(\ovk)}(X_1,\dots, X_n)$
    \item $(\alpha\cdot f)^{(\ovk)}(X_1,\dots, X_n)=\alpha\cdot f^{(\ovk)}(X_1,\dots, X_n)$.
  \end{enumerate}
\end{observation}
Since throughout this report we only use hasse derivatives in the context of univariate polynomials we will only focus on univariate polynomials and see some properties of hasse derivatives. Here we take $\bbF$ to be any field.
\begin{lemma}{}{}
  \begin{enumerate}[label=(\roman*)]
    \item For $f_1,\dots, f_t\in\bbF[X]$ then we have $$(f_1\cdot f_t)^{(n)}(X)=\sum_{\substack{n_1,\dots, n_t\geq 0\\ n_1+\cdots n_t=n}}\prod_{i=1}^tf_i^{(n_i)}(X)$$ 
    \item For any $\alpha\in\bbF$ we have $$\lt((X-c)^t\rt)^{(n)}=\binom{t}n\cdot (X-c)^{t-n}.$$  
    \item For $0\leq n\leq $ and $f,g\in \bbF[X]$we have $$\lt(gf^t\rt)^{(n)}(X)=\tdg\cdot f^{t-n}(X)$$ where $\tdg\in\bbF[X]$ with $\deg(\tdg)\leq \deg(g)+n(\deg(f)-1)$.  
  \end{enumerate}
\end{lemma}
\begin{proof}
  \begin{enumerate}[label=(\roman*)]
    \item Since the hasse derivative operator is linear it is enough to show this property for monomials. So suppose $f_j(X)=X^{k_j}$ for all $j\in[t]$. Then $f_1\cdots f_t=X^{\sum_{j=1}^t k_j}$. Therefore, $$(f_1\cdots f_t)^{(n)}(X)=\binom{\sum_{j=1}^t k_j}t,\qquad \prod_{j=1}^t f_j^{n_j}(X)=\prod_{j=1}^t \binom{k_j}{n_j}\ \text{ for all }n_j\in \bbZ_0, j\in[t]$$So it is enough to show that $$\binom{\sum_{j=1}^t k_j}t=\sum_{\substack{n_1,\dots, n_t\geq 0\\ n_1+\cdots n_t=n}}\prod_{j=1}^t \binom{k_j}{n_j}$$But this is true by comparing the coefficients of $X^n$ in $(X+1)^{\sum_{j=1}^t k_j}=\prod_{j=1}^t(X+1)^{k_j}$. So we have the result. 
    \item Now take $f_j(X)=X-c$ for all $j\in[t]$. Then $f_j^{(n)}(X)=1$ if $n=1$ and $0$ if $n>1$. Therefore, if we use the first part on $f_1,\dots, f_j$ in the sum on right-hand side the only terms which survives are when each $n_j$ are either $0$ or $1$. The number of such terms are $\binom{t}n$ and each term is $(X-c)^{t-n}$. 
    \item Again by using the first part we get $$(gf^t)^{(n)}(X)=\sum_{\substack{n_0,n_1,\dots, n_t\geq 0\\ n_0+n_1+\cdots +n_t=n}}g^{(n_0)}(X)\prod_{j=1}^t f^{(n_j)}(X)$$
  \end{enumerate}
\end{proof}
\begin{lemma}{}{lem:derivative-xq-x}
  Suppose $U(X), V(X) \in \bbF_q[X]$, $\deg(V) < q$, $i \geq 0$, and
  \[
    U(X) = V(X)\cdot \Lambda^i(X).
  \]
  Then for $0 \leq \ell < q$,
  \[
    U^{(\ell)}(X) \equiv (-1)^i \cdot V^{(\ell-i)}(X)
    \pmod{\Lambda(X)}.
  \]
\end{lemma}
\begin{Theorem}{}{}
  Suppose $\Char(\bbF)=p$ where $p$ is a prime. Let $Q(X,Y)\in\bbF[X,Y]$ be a polynomial and define $h(X)=Q\lt(X,X^{p^s}\rt)$ for some $s\in\bbN$. Then for all $0\leq n<p^s$ $$h^{(n)}(X)=Q^{(n,0)}\lt(X,X^{p^s}\rt)$$where $Q^{(i,0)}$ is the $i^{th}$ partial derivative of $Q$ with respect to $X$ where we consider $Q$ as a polynomial of $\bbF[Y][X]$. 
\end{Theorem}
\begin{proof}
  
\end{proof}